We choose \texttt{Qwen3-4B} as our smaller memory model because of its scale and long-context capability.
We always use the maximum summary length of $2048$ for consistency.
We pair the memory model with agent models with various sizes and capabilities.
Specifically, we choose \texttt{DeepSWE}, \texttt{Qwen3-Coder-Max} and \texttt{GLM-4.7}.
We compare with the following baselines: \textbf{Full-Context} is the standard agentic inference scenario with long-context inference. \textbf{No Summary} removes the summaries from the agent model prompt and leave the other components intact. \textbf{Qwen3-4B (base)} uses off-the-shelf model as memory model. Finally, $\textbf{GRPO}_{AC}$ is a model trained using our proposed framework in Section~\ref{sec:train}.
For \texttt{DeepSWE} and \texttt{Qwen3-Coder-Max}, we use $k=2$, and for \texttt{GLM-4.7} we use $k=4$.
For latency comparison, we start two different vLLM engines for both agent model and memory model. We consistently choose vLLM version \texttt{0.14.0} with prefix caching and chunked prefilling. We further compare latency under two scenarios: with or without CPU offloading, to demonstrate the capability of \model on wide deployment settings.
For \texttt{DeepSWE}, we run on A100 (80GB), and for the other two models, we run on H200.
Notice that in production, one memory model server can be used for multiple agent models. See further discussions in Section~\ref{sec:discuss}. Detailed hyperparameters for both training and evaluation are provided in Appendix~\ref{app:hyperparameters}.